\subsection{Finite support frames and the scalar ideal}
\label{subsec:finite-frame-ideal}

We prove Theorem~\ref{thm:finite-frame-ideal} and distinguish its scalar
ideal from the ordinary coefficient ideal. The finite dual-basis
argument is classical for semimodules \cite{DeorePatil}. For commutative
semirings, finite-projective dualizability is also described by
Niefield and Wood \cite[Corollary~2.2 and its proof, pp.~1223--1225]{NW}.
Niefield's quantale theorem \cite[Theorem~4.4, p.~848]{Niefield} relates
projectivity to total continuity in a category whose modules and maps
preserve arbitrary joins. Here only finite joins are required, and $Q$
need not be a quantale. We use no transfer of projectivity between these
two categories.

\begin{proof}[Proof of Theorem~\ref{thm:finite-frame-ideal}]
Choose a splitting $L\xrightarrow{s}Q^n\xrightarrow{r}L$ with $rs=1$.
Let $\lambda_i$ be the $i$th coordinate of $s$ and let $l_i=r(e_i)$.
Then $k=r(s(k))=\bigvee_i\lambda_i(k)l_i$, proving the frame identity.
In particular, each $\lambda_i$ is admissible at $l_i$.

Fix $l$. For a tuple $a_i\in K(l,l_i)$, the function
$\psi_a(k)=\bigvee_i\lambda_i(k)a_i$ is left $Q$-linear: right
multiplication by $a_i$ commutes with the left scalar action, and finite
joins of linear functions remain linear. Moreover,
\[
 \psi_a(k)l
 =\bigvee_i\lambda_i(k)(a_i l)
 \leq\bigvee_i\lambda_i(k)l_i=k.
\]
Thus $\psi_a\in D_l$. Conversely, if $\psi\in D_l$, put
$a_i=\psi(l_i)$. Admissibility gives $a_i l\leq l_i$, and applying
$\psi$ to the frame identity gives
\[
 \psi(k)=\bigvee_i\lambda_i(k)\psi(l_i)=\psi_a(k).
\]
This proves the parametrization without interchanging any scalar
factors.

Each $K(l,l_i)$ contains zero and is closed under finite joins. Their
finite product is a nonempty directed poset, and the map
$a\mapsto\psi_a$ is monotone and preserves finite joins. For any
$\psi\in D_l$, the poset of tuples with $\psi\leq\psi_a$ is nonempty
by the displayed representation. Coordinatewise joins give upper
bounds inside that same poset. Its upper comma categories are therefore
nonempty and directed, so the parametrization is final.

It follows in particular that
\[
 C(l,k)=\mathord\downarrow\{\psi(k):\psi\in D_l\}.
\]
The right side is intrinsic. Since $D_l$ is closed under pointwise
finite joins, its value set is nonempty and finite-join closed; its
downward closure has the same properties. Admissibility gives
$\psi(k)l\leq k$, proving $C(l,k)\subseteq K(l,k)$.

Evaluation $D_l\to C(l,k)$ is final as well. For each $q\in C(l,k)$
there is a functional with $q\leq\psi(k)$, and pointwise joins give
upper bounds for any two such functionals. Consequently
\[
 \varinjlim_{\psi\in D_l}F(\psi(k))
 \simeq\varinjlim_{q\in C(l,k)}F(q),
\]
compatibly with both maps into $\varinjlim_{K(l,k)}F$.
Theorem~\ref{thm:action-colimits} now gives \eqref{eq:frame-J}.
This is an equality of images, with no injectivity assertion.

If each scalar constraint set has a greatest element, its product has
the greatest tuple $(r(l,l_i))_i$. The resulting functional dominates
every $\psi_a$, proving \eqref{eq:frame-delta}. For finite $Q$, the
finite join of all elements of $K(l,l_i)$ supplies that maximum. With
the empty frame, $L$ is zero, the parameter product is a singleton, the
only functional is zero, and $C(l,k)=\{0\}$. All formulas and finality
arguments remain valid.
\end{proof}

The complete-join hypothesis in the quantale comparison cannot be
removed by taking ideal completions and then restricting a splitting.
Here $\operatorname{Idl}$ denotes downward-closed, finite-join-closed
subsets containing zero, with joins given by ideal closure of unions.

\begin{proposition}[ideal completion does not reflect projectivity]
\label{prop:completion-projectivity-boundary}
Let
\[
 Q=\{0\}\sqcup\{t^n:n\in\Z\},\qquad
 t^m+t^n=t^{\max(m,n)},\qquad t^m t^n=t^{m+n},
\]
where $0$ is the additive zero and multiplicative absorber. Give
$L=\B$ the action in which every nonzero scalar acts as the identity.
Then $L$ is cyclic and finitely presented but is not projective over
$Q$. Nevertheless $\operatorname{Idl}(L)$ is a retract of the rank-one
free complete module over the quantale $\operatorname{Idl}(Q)$.
\end{proposition}
\begin{proof}
The presentation with one generator $x$ and relation $tx=x$ identifies
all $t^n x$, by multiplication with positive and negative powers of $t$.
It gives exactly $L$, since the support map $Q\to\B$ preserves the
relation and distinguishes zero from the common nonzero class. If that
surjection had a $Q$-linear section, the image of $x$ would be $t^m$ for
some integer $m$, and linearity would imply
$t^{m+1}=t\,s(x)=s(tx)=s(x)=t^m$, a contradiction.

The ideals of $Q$ are precisely $\{0\}$, the principal ideals
$\mathord\downarrow t^n$, and $Q$ itself: a nonempty bounded-above
downward subset of integer exponents has a greatest element. Write
$\widehat Q=\operatorname{Idl}(Q)$ and $E=Q\in\widehat Q$. Its unit is
$\mathord\downarrow t^0$, whereas $E$ is an idempotent satisfying
$IE=E$ for every nonbottom ideal $I$ and $\{0\}E=\{0\}$. Indeed, any
nonzero $t^n\in I$ satisfies $t^n Q=Q$.

The completed module $\operatorname{Idl}(L)$ is again $\B$, with each
nonbottom ideal acting as the identity. The support map
$\widehat Q\to\B$ has the section $0\mapsto\{0\}$,
$1\mapsto E$. Both maps preserve arbitrary joins, and the displayed
product identities prove their $\widehat Q$-linearity. This is the
required rank-one retraction.

Finally, $E$ is the directed join of the proper principal ideals
$\mathord\downarrow t^n$ and is not compact. The completed section
therefore sends the compact nonzero element of $\B$ to a noncompact
ideal. It does not restrict to the original coefficient semiring.
Finite presentation of $L$ does not repair this failure.
\end{proof}

There is also a structural property of the intrinsic scalar sets.
Regard $K(l,k)$ as the hom set of a category with objects $L$, identity
$1_Q$ at every object, and composition $b\circ a=ba$. Its hom sets are
join-semilattices with zero; composition preserves finite joins. An
\emph{order ideal} of this category means downward-closed,
finite-join-closed subsets of its hom sets stable under composition on
both sides. Products of order ideals use downward finite-join closure
of composites. Arbitrary joins of scalar arrows are not assumed.

\begin{proposition}[idempotence of the scalar support ideal]
\label{prop:scalar-ideal-idempotent}
For every finitely generated projective support $L$, the sets $C(l,k)$
form a two-sided order ideal and satisfy $C^2=C$. Explicitly,
\[
 C(l,k)=\mathord\downarrow\left\{
  \bigvee_{j=1}^m b_j a_j:
  a_j\in C(l,h_j),\ b_j\in C(h_j,k),\ h_j\in L\right\},
\]
where $m\geq0$. This does not imply
$\mathcal J_L^2=\mathcal J_L$.
\end{proposition}
\begin{proof}
Take $a\in C(l,h)$ and $b\in K(h,k)$. Choose $\psi\in D_l$ with
$a\leq\psi(h)$. Linearity and monotonicity give
\[
 ba\leq b\psi(h)=\psi(bh)\leq\psi(k),
\]
so $ba\in C(l,k)$. For the other side, take $a\in K(l,h)$ and
$b\in C(h,k)$, and choose $\phi\in D_h$ with $b\leq\phi(k)$.
The function $\theta(x)=\phi(x)a$ is left linear and satisfies
$\theta(x)l=\phi(x)(al)\leq\phi(x)h\leq x$.
Thus $\theta\in D_l$ and $ba\leq\theta(k)$, proving the other ideal
condition. It follows that $C^2\subseteq C$.

Conversely, for $q\in C(l,k)$ choose $\psi\in D_l$ with
$q\leq\psi(k)$. The frame expands this value as
\[
 q\leq\psi(k)=\bigvee_i\lambda_i(k)\psi(l_i).
\]
Here $\psi(l_i)\in C(l,l_i)$. Since $\lambda_i$ is admissible at
$l_i$, we also have $\lambda_i(k)\in C(l_i,k)$. Each term is an
allowed composite for $C^2$, proving the reverse inclusion. Empty joins
cover the zero support.

The coefficient groups introduce a separate issue. A scalar inequality
$q\leq\bigvee_i b_i a_i$ need not express every coefficient from $F(q)$
as a sum of products from $F(b_i)$ and $F(a_i)$. No such generation or
surjectivity is among the coefficient-diagram axioms. More concretely,
Proposition~\ref{prop:canonical-arbitrary-ideal} realizes any two-sided
ordinary ideal $I$ as $\mathcal J_L$ at singleton support, while the
scalar ideal is $C=\{0\}$. Choosing $I=2\Z$ gives
$\mathcal J_L^2=4\Z\ne2\Z=\mathcal J_L$, although $C^2=C$.
\end{proof}
